\section{Reproduction and incorporation into ProveIt}
\label{sec:reproduce}

The packet is intended to be unpacked as a self-contained research
contribution. Its builders use Python's standard library, explicit local
paths, and pinned input hashes. They do not need network access or an
optimization search to reconstruct either delivered table. Assertions
must remain enabled: do not use Python's \code{-O} or \code{-OO} options.

From the packet root, the primary integrated replay is
\begin{verbatim}
python3 -B portable/reproduce.py \
  --full --output-dir /tmp/k222-replay
\end{verbatim}
Use an initially empty output directory outside the packet. The replay
checks input integrity; reconstructs the exact source; checks the new
closed logical proofs and emitter; reconstructs the literal parent,
folded, and canonically labelled tables; verifies the folds and startup;
and freshly rederives all primary runtime cases. The main README
documents optional predecessor and alternative replays and the exact
success labels. Retained runtime summaries are hash-bound to their
tables and sources, but they are distinguished from a fresh derivation.

The main artifact groups are:
\begin{center}
\begin{tabular}{>{\raggedright\arraybackslash}p{0.30\linewidth}>{\raggedright\arraybackslash}p{0.60\linewidth}}
\toprule
Path & Purpose\\
\midrule
\path{report/} & This report, editable LaTeX, generated complete table,
and PDF.\\
\path{portable/machine/} & Primary source, three labelings, exact parent,
maps, startup and validation data.\\
\path{portable/source/} & Deterministic source/emitter checker,
pinned macro inputs, new logical certificate, and theory note.\\
\path{portable/alternative/} & Conditional-K222 source, tables,
independent construction/checks, and receipts.\\
\path{portable/vendor/} & Unchanged pinned K224 archive with inherited
proofs, licenses and nested dependencies.\\
\path{research/} & Independent review and bounded exploratory results,
clearly separated from the positive proof.\\
\bottomrule
\end{tabular}
\end{center}

For repository integration, a suitable candidate directory is
\path{Computability/BusyBeaver/Research/ZFC/candidate-seed-g222/}.
The primary table, source, provenance, reproduction entry point,
theory note, and receipts should be kept together. The alternative can
be retained as an independent sibling or within the same research
packet. The report and complete packet can enter the existing intake
process. The root README and provenance identify which 223-state input
belongs to which construction, avoiding ambiguity between similarly
named incoming archives.

The table appendix is generated from the canonical file. Rebuilding the
PDF requires a standard LaTeX installation with the packages named in
\path{report/main.tex}. All files necessary to the delivered PDF are
included. Manifest hashes provide transport integrity; they do not
substitute for any mathematical proof.

\section{Other directions examined and their measured limits}

The successful search attacked both the theory and the encoding.
On the theory side, output-root K, conditional K, and seed-request K
were compared through their actual existing source programs, and the
merged equality theory was combined with the latter two. On the machine
side, source sharing, address constraints, decoder helper preservation,
and local state invariants were checked directly.

A small additional experiment changed $G$'s canonical binder from
$\cR(a,b)$ to $\cE(a,b)$ or $\cQ(a,b)$. The inherited general-helper
theorem permits these positive-tag builders: each is larger than both
endpoint labels. On the initial combined seed layout, the equality-code
binder produced 4,746 occurrences and the same 208 subtrees; the
quantifier-code binder produced 4,745 occurrences and 210 subtrees.
Neither yielded a smaller source in that experiment. The unrepaired
address patterns and full machine behavior of those variants were not
certified, and they are not additional claimed consistency machines.

A finite-context state-fold search also examined the three starting
223-state tables and an intermediate equivalent 222-state table. Its search class
used constant-bit invariants within radius eight and finite positive
simulation blocks of at most twenty steps. No additional one-state fold
was found in that class. This is an accurately scoped unsuccessful
search, not a lower bound. It does not exclude wider contexts,
different invariants, simultaneous multi-state transformations, or
different source layouts.
The retained receipts identify the exact searched bytes; the final
primary table's last direct-row simplification was not searched again.

More radical formalism changes remain possible but were not promoted
to state-count claims here. A direct standard-ZFC proof recognizer would
avoid the intermediate K interpretation but would need to implement
its logical and schema recognizers compactly. A finite class-theory
presentation could reduce axiom-schema machinery while increasing the
vocabulary and proof-checker overhead. A purely membership-based
source could eliminate primitive equality obligations while making
its equality translations and auxiliary names more expensive. None
of those qualitative observations supplies a literal table or a
certified state count. The successful calibrated K compositions are
the concrete results of this investigation.

\section{Further research questions}
\label{sec:future}

\begin{enumerate}
\item \textbf{Can the 208-subtree source retain thirteen aliases?}
The primary source has one fewer subtree than the conditional variant,
but the latter preserves one additional decoder alias. A compatible
208-subtree layout with thirteen aliases would suggest 221 states at
the same engine cost. The arithmetic is a target, not a compiler theorem:
an actual table must retain helper paths and pass full runtime checking.

\item \textbf{Can grammar selection optimize decoder sharing directly?}
The current construction measures repeated subtrees, then checks
whether contextual sharing survives. A joint solver could make source
parenthesization, protected helper paths, and state assignment part of
one optimization problem. A useful solver should emit independently
checkable exact source and table certificates, so its search procedure
does not enter the mathematical trust base.

\item \textbf{Can a stronger invariant remove a further engine state?}
The successful fold uses only two one-bit facts. Relations between
two tape cells, frame delimiters, or control phases might support a
reduction outside the searched radius-eight constant-fact class.
The old/new invariant distinction must be preserved when several
states take over each other's previously impossible branches.

\item \textbf{Can the seed request absorb another logical obligation?}
The $G$ construction succeeds because strengthening one row makes a
separate helper redundant and the original row can be recovered.
Other candidates should be tested first for literal interderivability
using the retained weak quantifier calculus, before source compression
is attempted. A semantic validity argument alone need not prove the
required native-calculus comparison.

\item \textbf{Can binder choice and arithmetic reuse be optimized together?}
The general helper theorem allows several safe numerical binder
functions. Merely replacing the formula-level binder did not improve
the tested subtree counts. Computing an equality atom once and reusing
its code, or exploiting a small tag difference between constructors,
may change the arithmetic source more substantially. Such a change
requires an arbitrary-stack macro contract and the same exact
matrix/private-code discipline.

\item \textbf{Can the new endpoint be made kernel-checked?}
The most direct formalization tasks are the uniform $G$ interderivation
for the exact seed bundle, the modified axiom-command interface, the
finite source certificate, and the 222-state table simulation. The
new all-tape folds are especially small candidates for an initial
Rocq proof. Completing those tasks would improve the verification
status without changing the machine.

\item \textbf{How weak a metatheory proves the entire equivalence?}
The present target is provability in ZFC. The new source transformations,
finite syntax translations, and affine/list inductions are elementary,
but weakening the global theorem to PA requires a precise internal
analysis of every inherited consistency reduction. A proof in an
external proof assistant does not by itself answer that question.
\end{enumerate}

\section{Attribution and conclusions}

Glazer's K theory, seed formulation, K0 quotient interpretation, and
consistency reductions are the mathematical foundation. The proof-search
and compiler lineage includes the work credited by the pinned packets:
Mira, Vladimir Reshetnikov's ProveIt project, Jeff's busy-beaver
development, Andrew James Wade, BusyCoq, and the Coq Library of
Undecidability. Existing licenses and notices remain in force for those
materials. The incoming seed223 and merged-equality223 packets supply
the two immediate ideas being combined.

This investigation supplies two explicit 222-state candidates and a
fully specified primary construction. The new source composition,
exact literal tables, uniform source binding, all-tape folds, complete
startup, and unbounded runtime checks form a reproducible improvement
over the inspected 223-state baseline. The proof chain and its inherited
premises are stated explicitly. Absolute consistency of ZFC,
unconditional non-halting, minimality, and a new end-to-end
proof-assistant endpoint are not consequences claimed by this work.
